\section*{Bab \ref{ch:b2:microbial-metabolism} --- Metabolisme Mikroba dan Biofilm}

\begin{solution}{exo:b2:microbial-metabolism:1}
Sianobakteri: foto-lito-autotrof (cahaya, air, $\mathrm{CO_2}$).
\emph{Nitrosomonas}: kemo-lito-autotrof (oksidasi amonium,
$\mathrm{CO_2}$). \emph{E.~coli} pada glukosa: kemo-organo-heterotrof.
Bakteri ungu nirbelerang pada suksinat dalam cahaya:
foto-organo-heterotrof. \emph{Thiobacillus} pada sulfida di dalam
gelap: kemo-lito-autotrof.
\end{solution}

\begin{solution}{exo:b2:microbial-metabolism:2}
Oksigen (beberapa milimeter permukaan), nitrat (tepat di bawahnya,
tempat oksigennya habis), oksida mangan dan besi (peralihan cokelat ke
kelabu), sulfat (lapisan sulfida hitam, tebal pada endapan laut),
$\mathrm{CO_2}$ bagi \omterm{def:b2:microbial-metabolism:anaerobic}{metanogenesis} (terdalam, atau tepat di bawah
permukaan pada air tawar yang miskin sulfat).
\end{solution}

\begin{solution}{exo:b2:microbial-metabolism:3}
$5\,\mathrm{C_6H_{12}O_6} + 24\,\mathrm{NO_3^-} + 24\,\mathrm{H^+}
\to 30\,\mathrm{CO_2} + 12\,\mathrm{N_2} + 42\,\mathrm{H_2O}$: 4,8
ion nitrat tiap glukosa (tiap nitrat menerima 5 elektron, tiap
glukosa memberi 24).
\end{solution}

\begin{solution}{exo:b2:microbial-metabolism:4}
Sebuah paguyuban mikroba yang melekat pada permukaan dan terbenam di
dalam matriks polisakarida, protein, dan DNA buatannya sendiri.
Tahapnya: pelekatan, mikrokoloni dan sekresi matriks, pematangan
bermenara dan berkanal, penyebaran. Plak dan karies gigi; jangkitan
kateter dan implan (serta paru-paru pada fibrosis kistik); pengotoran
dan pengaratan pipa serta lambung kapal --- atau, secara berguna,
saringan tetes pada pengolahan air.
\end{solution}

\begin{solution}{exo:b2:microbial-metabolism:5}
Delapan elektron jatuh dari \qty{-0.42}{V} ke \qty{-0.24}{V}:
$\Delta G^{\circ\prime} = -8\times 96.5\times 0.18 =
\qty{-139}{kJ/mol}$ tiap mol metana (nilai tabel \qty{-131}{kJ/mol}),
yakni \qty{35}{kJ} tiap $\mathrm{H_2}$: paling banyak 0,7 ATP tiap
hidrogen, dalam praktiknya kira-kira separuhnya. Energi tiap molekul
substratnya begitu kecil sehingga metanogen tumbuh lambat (\omterm{def:b2:unicellular-diversity:division}{waktu
generasi} berjam-jam sampai berhari-hari) dan mengubah hampir seluruh
substratnya menjadi gas dan bukan menjadi sel.
\end{solution}

\begin{solution}{exo:b2:microbial-metabolism:6}
$\mu = 1.2\,S/(5 + S)$: \qty{0.20}{h^{-1}} ($T = \ln 2/\mu =
\qty{3.5}{h}$), \qty{0.60}{h^{-1}} (\qty{1.2}{h}), \qty{1.09}{h^{-1}}
(\qty{0.64}{h}). Sembilan puluh persen: $S = 9K_s = \qty{45}{mg/L}$.
\end{solution}

\begin{solution}{exo:b2:microbial-metabolism:7}
$S^* = 0.05\times 0.4/(0.8 - 0.4) = \qty{0.05}{g/L}$; $X^* =
0.4\times(5 - 0.05) = \qty{1.98}{g/L}$; keluarannya $DX^* = 0.4\times
1.98 = \qty{0.79}{g/L/h}$; pencuciannya $D_c = 0.8\times 5/5.05 =
\qty{0.79}{h^{-1}}$.
\end{solution}

\begin{solution}{exo:b2:microbial-metabolism:8}
Yang tertangkap: $0.25\times 275 = \qty{69}{kJ}$, yakni 1,4 ATP tiap
amonium; $12/1.4 = 8.7$, kira-kira sembilan ion amonium tiap karbon ---
dan beberapa kali lebih banyak lagi begitu biaya pembuatan NADH dengan
mendorong elektron ke atas dari nitrit diperhitungkan. Seorang
heterotrof memperoleh karbonnya dalam keadaan sudah tereduksi dan
tertata, serta tiga puluh ATP tiap glukosa; seorang penitrifikasi harus
membeli energi sekaligus daya pereduksinya dari pemberi yang lemah dan
membangun tiap karbonnya dari $\mathrm{CO_2}$, sehingga ia mengubah
hanya sebagian kecil substrat yang diolahnya menjadi sel.
\end{solution}

\begin{solution}{exo:b2:microbial-metabolism:9}
$L_p = \sqrt{2\times 1.2\times 10^{-9}\times 0.2/0.5} =
\sqrt{9.6\times 10^{-10}} = \qty{31}{\micro m}$. Melipatduakan $c_0$
mengalikan $L_p$ dengan $\sqrt 2$: \qty{44}{\micro m}; melipatempatkan
$q$ memerohkannya: \qty{15}{\micro m}.
\end{solution}

\begin{solution}{exo:b2:microbial-metabolism:10}
Dari atas ke bawah: air dengan alga dan sianobakteri (fotosintesis
beroksigen); pita berkarat berisi pengoksidasi besi tempat
$\mathrm{Fe^{2+}}$ bertemu oksigen; selubung putih berisi
pengoksidasi belerang tak berwarna (\omterm{def:b2:microbial-metabolism:trophic}{kemolitotrof}) di tempat sulfida
bertemu oksigen; pita ungu berisi bakteri belerang ungu (\omterm{def:b2:microbial-metabolism:anoxygenic}{fotosintesis anoksigenik} pada
sulfida); pita hijau berisi bakteri belerang hijau (hal yang sama,
lebih tahan sulfida, lebih hemat cahaya); lumpur hitam berisi
pefermentasi dan pereduksi sulfat. (a) Bakteri belerang ungu
membutuhkan cahaya dari atas sekaligus sulfida dari bawah, sehingga
mereka duduk pada antarmuka tempat kedua landaiannya berpotongan.
(b) Pereduksi sulfat membuat $\mathrm{H_2S}$, yang mengendapkan besi
sebagai FeS hitam. (c) Tertutup bagi materi: belerang, karbon, dan
nitrogen berdaur antara bentuk teroksidasi dan tereduksi tanpa keluar;
terbuka bagi energi: cahaya masuk, bahang keluar, dan tanpa cahaya
setiap daurnya berhenti.
\end{solution}

\begin{solution}{exo:b2:microbial-metabolism:11}
Produksinya $pn$, kehilangannya $kA$: $A^* = pn/k$. $A_c = 10^{-8}\times
6.02\times 10^{23} = 6.0\times 10^{15}$ molekul tiap liter; $n_c =
kA_c/p = 5\times 10^{-4}\times 6.0\times 10^{15}/5 = 6\times 10^{11}$
tiap liter, yaitu $6\times 10^{8}$ tiap mililiter. Biakan padatnya
($10^{9}$) berada di atas ambangnya, air laut ($10^{5}$) empat orde
besar di bawahnya. Di dalam sebuah \omterm{def:b2:microbial-metabolism:biofilm}{biofilm} selnya duduk pada $10^{11}$
tiap mililiter matriks, dan matriksnya memperlambat lolosnya isyarat
(secara efektif menurunkan $k$), sehingga sebuah mikrokoloni berisi
beberapa ribu sel di dalam satu nanoliter mencapai kuorum sedangkan
sel yang sama yang tersebar dalam satu liter tidak akan pernah mencapainya.
\end{solution}

\begin{solution}{exo:b2:microbial-metabolism:12}
Fiksasi nitrogen (hanya \omterm{def:b2:unicellular-diversity:prokaryote}{prokariot} yang memiliki nitrogenase): jika
berhenti, nitrogen terpadu biosfernya akan terkuras habis oleh
\omterm{def:b2:microbial-metabolism:anaerobic}{denitrifikasi} dan produksinya akan runtuh dalam beberapa dasawarsa.
Nitrifikasi dan \omterm{def:b2:microbial-metabolism:anaerobic}{denitrifikasi} (hanya \omterm{def:b2:unicellular-diversity:prokaryote}{prokariot}): jika berhenti,
amonium akan menimbun dan nitrat lenyap, dan nitrogen tidak akan pernah
kembali ke udara. \omterm{def:b2:microbial-metabolism:anaerobic}{Metanogenesis} dan \omterm{def:b2:microbial-metabolism:anaerobic}{respirasi anaerob}
(hanya \omterm{def:b2:unicellular-diversity:prokaryote}{prokariot}): jika berhenti, bahan organik di dalam endapan tanpa
oksigen, rumen, dan rawa akan menimbun tanpa termineralkan dan karbon
akan meninggalkan daurnya; begitu pula, fotosintesis beroksigen adalah
penemuan bakteri, dan tanpanya tidak akan ada oksigen untuk
direspirasikan. Eukariot adalah tambahan yang terlambat dan sempit
secara kimia pada sebuah planet \omterm{def:b2:unicellular-diversity:prokaryote}{prokariot}.
\end{solution}

\begin{solution}{pb:b2:microbial-metabolism:1}
\textbf{1.} $-24\times 96.5\times (0.82 + 0.43) = \qty{-2900}{kJ/mol}$.
\textbf{2.} Nitrat: $-24\times 96.5\times 1.17 = \qty{-2700}{kJ/mol}$;
sulfat: $-24\times 96.5\times 0.21 = \qty{-490}{kJ/mol}$.
\textbf{3.} Oksigen menghasilkan lebih banyak, sehingga aerob
menyisihkan penitrifikasi balik dalam merebut bahan organiknya selama
oksigennya masih ada, dan oksigennya menekan reduktase nitratnya;
\omterm{def:b2:microbial-metabolism:anaerobic}{reduksi sulfat} membutuhkan ketiadaan oksigen \emph{dan} ketiadaan
nitrat, sehingga baunya berarti tangkinya kurang teraerasi dan
nitratnya habis --- lagi pula $\mathrm{H_2S}$ beracun dan mengikis
betonnya.
\textbf{4.} $0.4\times 2900/50 = 23$ ATP; $0.4\times 2700/50 = 22$;
$0.4\times 490/50 = 4$.
\textbf{5.} Kira-kira enam kali lebih banyak ($23/4$).
\textbf{6.} Dua ATP tiap glukosa berarti hasil kira-kira \num{0.1},
sehingga \qty{90}{\%} energi substratnya tetap tinggal di dalam
hasilnya (asam, alkohol, $\mathrm{H_2}$), yang diubah metanogennya
menjadi metana dengan perolehan yang sama kecilnya; sedikit energi,
sedikit biomassa, dan karbonnya pergi sebagai gas.
\textbf{7.} $10^4\times 0.3 = \qty{3000}{kg}$ COD tiap hari.
\textbf{8.} Biomassanya $0.4\times 3000 = \qty{1200}{kg}$ COD; sisanya
\qty{1800}{kg} COD teroksidasi dan membutuhkan \qty{1800}{kg}
$\mathrm{O_2}$.
\textbf{9.} $1800/2 = \qty{900}{kWh}$ tiap hari.
\textbf{10.} $10^4\times 0.04 = \qty{400}{kg}$ nitrogen;
$4.57\times 400 = \qty{1830}{kg}$ $\mathrm{O_2}$.
\textbf{11.} Pada keadaan tunak $\mu = 1/\theta \le \mu_{\max}$:
tercuci keluar di bawah $\theta = 1/\mu_{\max} = \qty{2}{d}$.
\textbf{12.} $D = \qty{0.1}{d^{-1}}$: $S^* = 1\times 0.1/(0.5 - 0.1)
= \qty{0.25}{g/m^3}$.
\textbf{13.} $S^* = 0.1/(0.25 - 0.1) = \qty{0.67}{g/m^3}$, dan batas
pencuciannya naik menjadi \qty{4}{d}; operatornya harus memperpanjang
umur lumpurnya (membuang lumpur lebih sedikit): pada
$\theta = \qty{20}{d}$, kembali diperoleh $S^* =
0.05/0.2 = \qty{0.25}{g/m^3}$.
\textbf{14.} $2.86\times 400 = \qty{1144}{kg}$ COD tiap hari.
\textbf{15.} Amonium efluennya $0.25\times 10^4 = \qty{2.5}{kg/d}$,
yaitu \qty{0.6}{\%} bebannya; sisanya, \qty{397}{kg} nitrogen, pergi
sebagai $\mathrm{N_2}$: \qty{99}{\%}.
\textbf{16.} Yang dihemat: \qty{1144}{kg} $\mathrm{O_2}$. Seluruh
kebutuhannya: $(1800 - 1144) + 1830 = \qty{2490}{kg}$ $\mathrm{O_2}$
tiap hari.
\textbf{17.} $1200\times 0.6\times 0.35 = \qty{252}{m^3}$ metana
tiap hari.
\textbf{18.} $252\times 36 = \qty{9070}{MJ} = \qty{2520}{kWh}$;
listriknya $0.35\times 2520 = \qty{880}{kWh}$ tiap hari.
\textbf{19.} Aerasinya $2490/2 = \qty{1245}{kWh}$ tiap hari: metananya
menutup $880/1245 = \qty{71}{\%}$ darinya.
\textbf{20.} $D_e\,c'' = q$, sehingga $c = c_0 + ax + qx^2/2D_e$;
dengan $c(L_p) = 0$ dan $c'(L_p) = 0$: $a = -qL_p/D_e$ dan $c =
(q/2D_e)(L_p - x)^2$, dengan $L_p^2 = 2D_ec_0/q$.
\textbf{21.} $L_p = \sqrt{2\times 1.5\times 10^{-9}\times 0.2/0.3} =
\sqrt{2\times 10^{-9}} = \qty{45}{\micro m}$.
\textbf{22.} $45/300 = \qty{15}{\%}$ bersifat aerob. Di bawahnya,
penitrifikasi balik mereduksi nitrat yang berdifusi turun dari lapisan
permukaan yang menitrifikasi menjadi $\mathrm{N_2}$, dengan memakai
bahan organik yang berdifusi masuk.
\textbf{23.} Di dalam satu lapisan, landaian oksigennya menciptakan
lapisan aerob (nitrifikasi) di atas lapisan tanpa oksigen
(\omterm{def:b2:microbial-metabolism:anaerobic}{denitrifikasi}), berjarak beberapa puluh mikrometer; tangki
beraerasinya beroksigen di seluruh bagiannya, sehingga
\omterm{def:b2:microbial-metabolism:anaerobic}{denitrifikasinya} membutuhkan tangki tanpa oksigen tersendiri.
\textbf{24.} Klorinnya bereaksi dengan dan dihabiskan oleh polimer
matriksnya sebelum mencapai sel yang dalam; sel yang dalam itu
kelaparan, tumbuh lambat, atau dorman dan jauh kurang peka; hanya
lapisan permukaannya yang mati dan lapisannya tumbuh kembali dari bawah.
\textbf{25.} Kebutuhan oksigennya \qty{2490}{kg} tiap hari; metananya
\qty{252}{m^3} tiap hari; metananya memasok \qty{71}{\%} listrik
aerasinya (\qty{880}{kWh} dari \qty{1245}{kWh}).
\end{solution}
